\section{Meditation on a World Conception}

\begin{quotex}
Let no one ignorant of geometry enter. \flright{Attributed to \textsc{Plato}}

A little philosophy inclineth man's mind to atheism, but depth in philosophy bringeth men's minds about to religion. \flright{\textsc{Francis Bacon}}

Thales the Milesian held that self-knowledge is the beginning of every virtue … The man who knows himself knows that he is a rational, spiritual being, possessing the power of free choice and self-control, and hence religious, social, ethical, and in his spiritual nature, immortal. \flright{\textsc{St. Nectarios}, Metropolitan of Pentapolis}

\end{quotex}
In his \emph{Life of Dante}, \textbf{Giovanni Boccaccio} explained that truths may be hidden in poetry because there is value in having to make the effort to extract them. For the same reason, we read difficult works in metaphysics because working through the ideas and arguments makes them alive to us. That is why the Cliff Notes version of a philosophy, which reduces complex thoughts to a few key points, sound so shallow. Seriously, can someone learn mathematics simply by looking up the answers in the back of the book?

Let's look at another formulation of Plato's three degrees of knowledge:

\begin{itemize}
\item First Principles (episteme) 
\item Likely Stories (dianoia) 
\item Opinions (doxa) 
\end{itemize}
Opinions refers to the mass of facts around us and ideas that pop into our minds. There are large areas in life where it is possible to live like that. Nevertheless, the more intelligent will try to organize the mass of facts by inventing a likely story, à la Sherlock Holmes, to account for all those facts.

Scientific theories are a more sophisticated version of the likely story. The scientist will develop a theory, mathematical formula, etc., to ``save the appearances", i.e., to account for observed phenomena. That theory may come to be regarded as a ``law of nature", objectively true, apart from human subjectivity.

However, are those laws really discovered, or are they created? For example, what is the epistemological status of Newton's theory of gravity? His system is an intellectual tour de force. In a simple formula, he related disparate phenomena such as an apple falling from a tree to the orbits of planets around the sun. It even predicted unobserved events such as a feather and a stone falling at the same rate in a vacuum. But it did not explain gravity itself.

For quite some time, educated people regarded it as an ineluctable law of nature, at least until Einstein's General Theory of Relativity overturned it. At that point, the Principia could not be regarded as revealing the ultimate secrets of gravity. Hence, it can only be a brilliant creative effort from Newton, not an objective description of reality.

Nevertheless, educated, and not so educated, people remain obsessed with the likely story. A prime obstacle to communication is that religion, myths, etc., are regarded as crude likely stories, invented to explain facts, but are now superseded by science.

This is the bottom up approach: start with the evidence of the senses and use reason to formulate a likely story. Of course, there is a pecking order to likely stories. The better ones will explain a wider range of phenomena, while the best will predict phenomena not yet observed. This is never ending and never provides certainty, since a better story may be created next week.

The other approach is to start from the top, in the manner of geometry. There are postulates and axioms, without which, no geometry is even possible. We don't do geometry by observing triangles in nature and then formulate theories about them. Rather, the properties of triangles are deduced from the first principles; any triangles existing in nature can only be better or worse approximations of ideal triangles.

Valentin Tomberg offers us some first principles, not a theory, from which to understand the world. Their truth is revealed by meditating on them, not by blind faith, nor by doing experiments. What follows is a personal meditation, not a proof; you can try your own at home. The following three points are the foundation.

\begin{itemize}
\item The world is an ordered whole. 
\item The world order is of a moral nature. 
\item Mankind is called to arrange its own life in harmony with the moral world order. 
\end{itemize}
The first thing to note is the simplicity of these three points. They do not require an advanced knowledge of philosophy or science to be understood. People uninterested in those endeavours can accept them, although they may need to be expressed in religious or symbolic languages.

The second thing to note is that they are universal, i.e., they should be acceptable to any of the valid Traditional religions. For example, Tao is the principle of order in Taoism, rta in Hinduism, Logos in Christianity.

\paragraph{Ordered Whole}
The world forms a unity, and everything is ultimately related, both in time and space. Even in theories of the multiverse, there are common physical laws that account for them. That is the higher unity.

Moreover, there is an order to the world. Things and events are causally related to each other. Pure chance explains nothing. Causation is not necessarily on the same level. If a visible effect seems to have no visible cause, then the cause can be traced to a higher level.

Philosophical systems like those of Spinoza and F H Bradley develop this conceptual point. However, they envision a sort of block universe which omits and denies the next point.

\paragraph{Moral Nature}
Moral convictions such as rights, duties, justice, etc., insert a new element into the block universe. They can only make sense assuming free will, i.e., the consciousness of making choices.

Such a world requires a field in which it is possible to make choices, so there must be a not-I to the I. That is, there are necessarily objects or people which can obstruct my will.

It also presupposes a knowledge of good and evil. A world with free will is not predictable by any science.

\paragraph{Moral World Order}
If minerals, plants, and animals follow laws by necessity, there are moral laws — which transcend animal existence — that may or may not be followed. One must prefer good over evil, but to prefer one thing over another is to love it.

To overcome obstacles to the good requires creativity. Hence love and creativity are introduced into the world. A world of moving atoms has no need for those qualities.

Now it is clear that such a world must necessarily exist, given the preference for the good, the existence of love, and creativity. It cannot arise by accident from a ``Big Bang" or random variations. This is because:

\begin{itemize}
\item Being is preferable to non-being 
\item Order is preferable to chaos 
\item Beings with free will are preferable to unconscious objects 
\end{itemize}
In summary, the world follows from these fundamental points; they are not part of a likely story that is formulated from observing the world.

\paragraph{Objections}
It can be helpful in your meditations to consider a world that denies any one or more of these points. Consider points like these:

\begin{itemize}
\item Can such a world actually exist? 
\item Can I live effectively and consistently in such a world? 
\end{itemize}
We have examples of laws that serve as guides for living in harmony with the moral world order. Examples are those of great lawgivers like Solon, Manu, Zoroaster, Moses, and so on.

Nihilists will deny most or all of the three points. Then they need to create an alternate world conception to inhabit. You make your own choice.



\flrightit{Posted on 2018-09-20 by Cologero }

\begin{center}* * *\end{center}

\begin{footnotesize}\begin{sffamily}



\texttt{Mikkel on 2018-09-23 at 15:22 said: }

``Seriously, can someone learn mathematics simply by looking up the answers in the back of the book? "

No, but I would like it if metaphysical books had some solutions in the back when I'm working on problems to check my efforts, now that I think of it, they probably do, I need to read harder, or perhaps the tip is personal meditation which you suggested, I will try this.

``The other approach is to start from the top, in the manner of geometry."

Man and His Becoming (pg. 34): ``To make a comparison with mathematics by way of clarification, it is thus that the geometrical point is quantitatively nil and does not occupy any space, though it is the principle by which space in its entirety is produced, since space is but the development of its intrinsic virtualities."

So then we must start at these intrinsic virtualities to understand, or at least we have to understand those intrinsic virtualities of a thing to truly understand/know it (not necessarily do we always have the benefit of being able to start clearly at that place, maybe too temporal a way of looking at things).

``One must prefer good over evil, but to prefer one thing over another is to love it.

To overcome obstacles to the good requires creativity." 

I think to love is to want the good to be expressed, ie. To love another is to want that good to be expressed in that individual, for that individual to be expressive of the true Self, which must be good and moral (to want that for another one must also know another). This reminds me of Boethius, in that all seek the good in every action that they do, however ignorance hinders these efforts greatly. Some seek to reduce the amount of actions they take, but elimination of that ignorance is ultimately what matters. Lately many resources I read outside of this `circle' are explaining the work of creativity and expression and their position to humanity, their arguments have an inherent assumption that these actions are good without ever stating why. The missing conclusion however, is that the good must be the end goal of these actions and should be acknowledged as such, but these actions occur from principle where it started from. Is this the from same idea as, ``that which is first or greatest in the principial order is … last and smallest in the order of manifestation."

A question to this however, what aspect about creativity specifically that overcomes obstacles? or maybe…why do obstacles require creativity to be passed? Reminds me of the The Augean Stables, not any man can do it.


\end{sffamily}\end{footnotesize}
